\section{Formal symmetries and coherent transport}
\label{sym:section}

We apply the affine-field construction of Theorem~\ref{hull:main}.
Its full ambient jets, and their holonomicity on both axes, are essential
here.  Smoothness of the values on an axis alone would not suffice.
We construct the formal strict contact symmetries of these fields,
prove that their projected values span the section tangent space,
and derive the entropy trace identity used in the next section.
All vector spaces in this section are complexified unless a real
structure is specified.  Formal series carry the coefficientwise linear
topology: a neighborhood of zero prescribes the vanishing of finitely
many Taylor coefficients.  In particular, no convergence of a formal
series at a nonzero point is asserted.

\subsection{The symmetry equations}

Fix a regular center and a section through it, with coordinates
$q\in\R^{2m}$, center $q=0$, and contact lift
\[
 \alpha=d\tau+\theta,\qquad \omega=d\theta.
\]
A strict contact vector field commuting with $X=\partial_\tau$ is
uniquely specified by a formal function $p(q)$:
\begin{equation}
 \label{sym:hamiltonian-lift}
 \iota_{Y_p}\omega=-dp,\qquad
 Z_p=Y_p+\bigl(p-\theta(Y_p)\bigr)\partial_\tau,
 \qquad \alpha(Z_p)=p.
\end{equation}
Indeed Cartan's formula gives $\mathcal L_{Z_p}\alpha=0$, and every
strict lift has this form. For formal functions $p,r$, we use the bracket
\[
 \{p,r\}=Y_pr=\omega(Y_p,Y_r),
 \qquad [Z_p,Z_r]=Z_{\{p,r\}}.
\]

Augment each affine field $\mathcal A^\varepsilon$ by a scalar slot:
\[
 \mathcal B^\varepsilon
   =\Span\{(a,1):a\in\mathcal A^\varepsilon\},
 \qquad \varepsilon\in\{u,s\}.
\]
Thus the scalar-one slice recovers the affine field, including its
translation part.  On the section, put $b_p=p-\theta(Y_p)$.
The infinitesimal action on an augmented section is
\begin{align}
 \mathcal D_p^u(\beta,s)
     &=(\mathcal L_{Y_p}\beta-hb_p\beta,\,Y_ps),
       \label{sym:augmented-action-u}\\
 \mathcal D_p^s(\beta,s)
     &=(\mathcal L_{Y_p}\beta+hb_p\beta,\,Y_ps).
       \label{sym:augmented-action-s}
\end{align}
These formulas retain the factors $e^{-h\tau}$ and $e^{h\tau}$ in the
contact lift.  Define $\mathfrak h$ by requiring each operator to
preserve its corresponding formal subbundle.  In a Grassmann chart,
choose a frame matrix $B^\varepsilon$ and a quotient matrix
$Q^\varepsilon$ with kernel $\mathcal B^\varepsilon$; the equations are
\begin{equation}
 \label{sym:linear-equations}
 Q^\varepsilon\mathcal D_p^\varepsilon B^\varepsilon=0,
 \qquad \varepsilon=u,s.
\end{equation}
They are linear in $p$.  Their coefficients through order $r$ use only
finitely many background coefficients and the $(r+2)$-jet of $p$.
Changing a frame or a Grassmann chart gives equivalent equations.
Consequently $\mathfrak h$ is a closed linear subspace of the formal
Hamiltonians and is closed under the Poisson bracket.  In particular it
is a linearly compact Lie algebra with continuous bracket.  Set
\begin{equation}
 \label{sym:isotropy-functional}
 \xi(p)=p(0),\qquad
 \mathfrak k=\{p\in\mathfrak h:Y_p(0)=0\}.
\end{equation}
Both the functional and the evaluation map are continuous;
$\mathfrak k$ is an open subalgebra.

\subsection{A signed dilation}

The normal forms used to construct the hulls had one-signed state
weights.  For the present purpose we need a different normalization,
of the section map itself, with both dynamical signs.  This is only a
formal normalization, performed one ordinary Taylor degree at a time.

\begin{lemma}[Signed symplectic normalization]
\label{sym:graded-contact}
Along almost every orbit of $f=\phi_T$ there are centered formal
symplectic section coordinates, with tempered jets at every fixed
order, having the following properties.  The symplectic form is a
constant form $\omega_0$.  If $\lambda_1,\ldots,\lambda_{2m}$ are the
signed Lyapunov exponents, counted with multiplicity and measured per
iterate of $f$, then
\[
 D=\sum_i\lambda_iq_i\partial_{q_i}
\]
is symplectic and every normalized section map commutes with $D$.
The coordinates lift strictly with
$\theta_0=\tfrac12\iota_R\omega_0$, where
$R=\sum_iq_i\partial_{q_i}$, and center height zero.  The centered return
map has the form
\begin{equation}
 \label{sym:graded-return}
 (q,\tau)\longmapsto(F(q),\tau+r(q)),
 \qquad r(0)=0,\quad Dr=0.
\end{equation}
The actual strong stable and unstable Taylor germs are respectively
the negative and positive coordinate axes, with their horizontal
lifts.
\end{lemma}

\begin{proof}
The symplectic pairing between Lyapunov spaces of exponents
$\lambda,\mu$ vanishes unless $\lambda+\mu=0$.  To see this, transport
the pairing along the orbit and use the two-sided Lyapunov estimates;
when the sum is nonzero, transport in the direction in which the
product of the two norms tends to zero.  A measurable symplectic
Oseledets frame therefore pairs opposite spaces and has tempered
finite-dimensional change-of-frame costs.  Centered Darboux
coordinates with these first derivatives give the initial constant
symplectic form.  The field $D$ preserves $\omega_0$.

Suppose that the section map has already been made weight-preserving
through degree $j-1$, with $j\ge2$.  Factor out its linear part.  The
non-weight-preserving part of its degree-$j$ term is a homogeneous
symplectic vector field.  More explicitly, the symplectic pullback
identity in degree $j-1$ has a linear term
$\mathcal L_V\omega_0$ and terms determined by lower degrees.  Those
lower terms have weight zero; its nonzero-weight part is therefore
$\mathcal L_{V_{\ne0}}\omega_0=0$.  A homogeneous symplectic vector
field is Hamiltonian, by the polynomial Poincare lemma.

On each nonzero-weight component solve the nonresonant homogeneous
coordinate-change equation, and take zero correction on the
zero-weight component.  Explicitly, write the degree-$j$ error after
factoring the linear map as $V_n$, and the coordinate correction as
$Q_n$.  Conjugating by the correction gives the equation
\[
 Q_{n+1}=A_{n*}Q_n-A_{n*}V_n,
 \qquad (A_{n*}Q)(q)=A_nQ(A_n^{-1}q).
\]
The operator $A_{n*}$ preserves Hamiltonian vector fields and has
the single exponent $-\gamma$ on their $D$-weight-$\gamma$ module.
Thus Lemma~\ref{nf:homological} applies for $\gamma\ne0$.  At this
degree the nonzero exponent differences form a finite set bounded
away from zero.  The appropriate forward or backward Green series
consequently gives a tempered correction.
The formal time-one map of its Hamiltonian vector field is
symplectic and changes no lower degree.  Induction gives compatible
coordinates at all orders.  This argument imposes no discreteness
condition on the union of the weights over all degrees.

The radial primitive $\theta_0$ satisfies
$\mathcal L_D\theta_0=0$.  A symplectic coordinate change has a strict
lift because the difference of its pulled-back primitive and the
original primitive is formally exact.  Normalize its primitive at
zero.  For the return map, strictness gives
$dr=\theta_0-F^*\theta_0$.  Both terms are $D$-invariant, so $d(Dr)=0$;
evaluation at zero gives $Dr=0$.

Finally, a graded section map preserves each signed linear axis:
a monomial involving only negative-weight inputs cannot have positive
output weight, and conversely.  The actual stable germ is a graph
over the negative axis with tempered finite jets.  If the first
nonzero graph coefficient had degree $j$, its invariance would give
a homogeneous recurrence whose exponent is a positive output weight
minus a sum of negative input weights.  This exponent is nonzero.
A homogeneous recurrence of nonzero exponent has no nonzero
two-sided tempered solution.  Thus this coefficient vanishes;
induction removes every graph coefficient.  The unstable argument
uses the opposite signs.  The strong leaves are horizontal.  On
either signed linear axis $\theta_0$ vanishes, so their centered
horizontal lifts have height zero as formal germs.
\end{proof}

\begin{proposition}[Dilation in the symmetry algebra]
\label{sym:dilation}
In the coordinates of Lemma~\ref{sym:graded-contact},
\begin{equation}
 \label{sym:dilation-hamiltonian}
 H=-T+\theta_0(D),\qquad Z_H=D-T\partial_\tau
\end{equation}
belongs to $\mathfrak k\subset\mathfrak h$.
The operator $\ad H$ on formal Hamiltonians is $D$ and is semisimple
on every finite Taylor quotient, with real weights.  Every closed
$\ad H$-invariant linear subspace admits the continuous projections
onto each individual weight.  On a finite-dimensional continuous
invariant subquotient, weight-zero vectors have weight-zero lifts.
Moreover $\mathfrak h_0\subset\mathfrak k$, and
$\xi$ vanishes on every nonzero weight space.
\end{proposition}

\begin{proof}
For clarity, consider the finite-dimensional space of augmented
section jets through a fixed order $N$.  Its subspace consisting of
jets valued in $\mathcal B^u$ is measurable and invariant under the
return cocycle.  In the normalized charts that cocycle acts on the
form component by
\begin{equation}
 \label{sym:entropy-jet-action}
 \beta\longmapsto e^{hT-hr}F^*\beta,
\end{equation}
and on the scalar component by composition with $F$.  For
$\mathcal B^s$ the multiplier in
\eqref{sym:entropy-jet-action} is $e^{-hT+hr}$.
These are pullbacks from the next center to the present center; they
may equivalently be viewed as a cocycle over the backward base map.

For an $m$-form monomial
$q^a\,dq_{i_1}\wedge\cdots\wedge dq_{i_m}$, its ordinary weight is
\[
 \sum_i a_i\lambda_i+\lambda_{i_1}+\cdots+\lambda_{i_m}.
\]
Add $hT$ for the $u$ form component and subtract $hT$ for the $s$
form component.  Scalar monomials have their unshifted ordinary
weight.  The return cocycle preserves these augmented weights:
$F$ is graded and $r$ has weight zero.  Each block has the single
Lyapunov exponent indicated by its weight.  Indeed its diagonal
ordinary-degree blocks are tensor representations of the linear
return maps with precisely that exponent.  The remaining terms are
triangular in ordinary degree; at a fixed truncation there are only
finitely many such extensions.  Their coefficients are tempered,
and the finite iterated sums in an upper triangular product add
only subexponential costs to the common exponent.

Lemma~\ref{nf:blocks} now splits the measurable invariant jet subspace
into these weight blocks.  The ambient finite jet cocycle of the
smooth return maps on the compact phase space has bounded one-step
coefficients, and its formal change of coordinates is tempered at
the fixed order.  No tempered bound on the hull subspace itself is
being assumed.
The infinitesimal diagonal action is
$(\mathcal L_D+hT, D)$ on the two augmented $u$ slots and
$(\mathcal L_D-hT,D)$ on the $s$ slots.  It therefore preserves both
augmented subbundles through order $N$.  Since $N$ was arbitrary,
Equations~\eqref{sym:augmented-action-u}--\eqref{sym:augmented-action-s}
show that $D-T\partial_\tau$ preserves both formal affine fields.
The identity
$d(\theta_0(D))=-\iota_D\omega_0$ proves
\eqref{sym:dilation-hamiltonian}.

For any Hamiltonian $p$, $\{H,p\}=Dp$.  Thus the asserted
diagonalization is literal monomial diagonalization.  To project
onto a prescribed weight $\gamma$, choose on each finite Taylor
quotient the polynomial spectral projector onto $\gamma$ for $D$.
Applied to a fixed formal series these projectors eventually agree
at every fixed coefficient, and converge coefficientwise to its
weight-$\gamma$ part.  A closed invariant subspace contains this
limit.  A continuous finite-dimensional quotient is detected at a
finite Taylor order, so the same projectors commute with passage to
that quotient and give the claimed lifts.  The argument applies also
to closed invariant subspaces before taking such a quotient.

All coordinate weights in degree one are nonzero.  A weight-zero
Hamiltonian therefore has zero differential at zero and belongs to
$\mathfrak k$.  Evaluation at zero sees only the constant monomial,
which has weight zero.
\end{proof}

\begin{lemma}[Entropy trace identity]
\label{sym:trace}
For every $p\in\mathfrak h_0$, the linear map
$A_p=dY_p(0)$ preserves both actual signed tangent planes and
\begin{equation}
 \label{sym:geometric-trace}
 h\xi(p)=\tr\bigl(A_p|_{E^s}\bigr).
\end{equation}
Here the strong tangent planes are identified with their projections
to the section.
\end{lemma}

\begin{proof}
Since $[D,Y_p]=0$ and $Y_p(0)=0$, the derivative $A_p$ commutes with
$D$.  It preserves the signed planes, which are the actual strong
planes by Lemma~\ref{sym:graded-contact}.  At the center, the actual
form $u$ annihilates $E^u$ and is a volume form on $E^s$.  Its
infinitesimal transform under $Z_p$ is therefore
\[
 \bigl(\tr(A_p|_{E^s})-h\xi(p)\bigr)u.
\]
Preservation of the affine field says that this vector belongs to
the direction space of $\mathcal A^u(0)$.  Pairing any such direction
with the actual opposite form $w(0)$ gives zero, by the constant
cross-pairing property in Theorem~\ref{hull:main}.  Since
$u(0)\wedge w(0)\ne0$, the displayed scalar is zero.  Complexification
causes no change in this argument.
\end{proof}

\subsection{Finite solution spaces on a whole axis}

The dilation coordinates above are only measurable in the center.
We will never differentiate them with respect to the center.
Instead we use the equations in fixed smooth coordinates and the
following finite-dimensional consequence of whole-axis coherence.

\begin{lemma}[Full-solution projection bundles]
\label{sym:full-projections}
Refine the regular set by a countable conull intersection.  For every
finite order $r$, the projections of $\mathfrak h$ and
$\mathfrak k$ to the Hamiltonian $r$-jets have smooth extensions on
each entire strong axis through a regular center.  Their fibers at
every point of that axis are the projections of the full formal
solution spaces for the extended background jets there.  Projections
between successive jet orders are surjective smooth bundle maps.
These statements also hold after appending a fixed finite list of
linear jet conditions whose matrices are smooth along the axes and
equivariant under the return maps.
\end{lemma}

\begin{proof}
Use the background extensions of Lemmas~\ref{hull:coherence} and
\ref{hull:measurability}.  Write $K_N(a)$ for the finite space of
$(N+2)$-jets satisfying Equation~\eqref{sym:linear-equations} through
degree $N$ at center $a$.  Isotropy, when required, is an additional
linear condition.  For $N+2\ge r$, set
\[
 V_{r,N}(a)=\pi_rK_N(a).
\]
Both $K_N$ and the map defining this image are finite-dimensional
linear algebra in finitely many background jets.  Their matrices
are smooth on the axis.

The ranks of the defining matrix $A_N$ and the augmented matrix
obtained by stacking $A_N$ with $\pi_r$ are invariant under the
return dynamics, because finite symmetry conditions and jet
projection are natural under centered coordinate changes.  By
ergodicity these ranks are constant almost everywhere,
simultaneously for all $r,N$ under consideration.  These two ranks
determine $\dim K_N$ and $\dim V_{r,N}$.  At a regular center a
nonzero maximal minor remains nonzero on a small axis patch.  Dense
regular agreement prevents a larger rank anywhere on that patch:
otherwise an open set of larger rank would meet the dense regular
set.  Thus both ranks, and consequently the kernel and image ranks,
are constant near the center.  Constant-rank linear algebra gives
smooth kernel and image bundles there.

This local assertion extends to the entire typical axis.  Fix the
finite list of matrices just used.  Their local patches have
measurable positive radii after also recording finite coefficient
bounds and a lower bound for the chosen nonzero minors.  Such data
are provided by Lemma~\ref{hull:measurability}.  Some positive-measure
set has these radii bounded below.  Returns of a typical center to
that set, in the direction contracting the axis, bring every fixed
compact portion of the axis into one of these patches.  Pullback
preserves the finite equations and both ranks.  This proves the
constant-rank conclusion on that compact portion.  Exhausting the
axis and intersecting over the countable list of finite tests gives
the assertion for the whole axis.  Forward time is used for the
stable axis and backward time for the unstable axis.

For fixed $r$, the spaces $V_{r,N}(a)$ decrease with $N$ inside a
finite-dimensional vector space.  At regular centers all their
dimensions are the same deterministic sequence of integers.  Choose
$N(r)$ after its eventual stabilization; a first adjacent equality
would not suffice.  The preceding whole-axis rank argument applies
to every $N$, so for every point of that axis
\begin{equation}
 \label{sym:stabilized-projection}
 V_r(a):=\bigcap_N V_{r,N}(a)=V_{r,N(r)}(a).
\end{equation}
It remains to identify this intersection with full solutions.

We use the elementary linear compactness principle: a system of
linear equations in coefficient variables, each equation involving
only finitely many variables, has a solution if every finite
subsystem has a solution.  One proof regards each equation as a
finitely supported row vector.  Finite consistency makes the
specified right sides a well-defined linear functional on the span
of the rows; extend it to the direct sum of all coefficient
coordinates.  Its values on the coordinate vectors give a solution.
If $v\in V_r(a)$, impose the prescribed $r$-jet $v$ together with
all formal symmetry equations.  Every finite subsystem is solvable
by the defining intersection, so this principle gives a full
solution with that jet.  Conversely every full solution projects
into each $V_{r,N}$.  This proves the claimed equality.

Equation~\eqref{sym:stabilized-projection} makes $V_r$ a smooth
bundle.  A full solution lifting any of its elements also supplies
a lift at order $r+1$, so $V_{r+1}\to V_r$ is surjective in every
fiber.  It is a smooth map of constant-rank bundles.  All arguments
apply unchanged to the isotropy condition or a finite list of
linear conditions with the stated smoothness and equivariance.
\end{proof}

\begin{proposition}[Transitivity]
\label{sym:transitivity}
At every regular center the map
\[
 \mathfrak h\longrightarrow T_0\R^{2m},\qquad p\longmapsto Y_p(0)
\]
is onto.  It remains onto at every point of the smooth formal
extensions on either entire typical axis.
\end{proposition}

\begin{proof}
Fix a regular center $a(0)=0$, and a smooth path $a(t)$ in one of its
actual axes.  Fix a desired equation order $r$.  By
Lemma~\ref{sym:full-projections}, extend the $(r+3)$-jet of its
particular isotropy $H$ to a smooth section $p_t$ of the full isotropy
projection bundle on that axis.  Each $p_t$ is represented by a
Taylor polynomial in $z=q-a(t)$.  It satisfies the symmetry equations
through degree $r+1$ and $Y_{p_t}(a(t))=0$.

Differentiate while holding the original coordinate $q$ fixed.
For a moving-center series $C_t(z)$ this differentiation is
\[
 \partial_t^{\mathrm{fixed}}C_t
   =\partial_t C_t-\sum_i a_i'(t)\partial_{z_i}C_t.
\]
Full ambient holonomicity of the background states that its
fixed-coordinate derivative vanishes, coefficient by coefficient.
The same is true for smooth frames and quotient matrices obtained
from it by the Grassmann-chart operations.  Differentiating the
linear equations through degree $r+1$ therefore shows that
$\partial_t^{\mathrm{fixed}}p_t$ satisfies them through degree $r$.
The extra degree is what compensates for differentiation of the
moving Taylor center; the order-two Hamiltonian operator accounts
for the choice of the $(r+3)$-jet.

Differentiating isotropy by the ordinary chain rule gives
\begin{equation}
 \label{sym:differentiated-value}
 Y_{\partial_t^{\mathrm{fixed}}p_t}(0)
   =-dY_H(0)a'(0).
\end{equation}
At the chosen center, $dY_H(0)$ is conjugate to the signed dilation
$D$. It is invertible and preserves each actual signed tangent plane.
Using paths in the two axes therefore gives, at every prescribed
finite equation order, a solution with any prescribed tangent
value.  Finite sums handle a value with both signs.  Apply the
linear compactness principle from
Lemma~\ref{sym:full-projections}, with this tangent value prescribed,
to obtain a full formal solution.  No compatibility of the
separately chosen finite families is needed.

Finally, the dimension of the evaluation image is the difference
between the dimensions of the full symmetry and isotropy
projections at first jet order.  Those dimensions have the same
constant values on the entire typical axes by
Lemma~\ref{sym:full-projections}.  Surjectivity at their regular
points consequently implies surjectivity throughout their smooth
extensions.
\end{proof}

\begin{proposition}[The effective symplectic homogeneous algebra]
\label{sym:effective}
The open isotropy $\mathfrak k$ contains no nonzero ideal of
$\mathfrak h$, and is exactly the coadjoint stabilizer of $\xi$.
Evaluation identifies
\[
 E=\mathfrak h/\mathfrak k\simeq T_0\R^{2m},\qquad
 \omega_E(\bar p,\bar q)=\xi(\{p,q\}).
\]
This form is nondegenerate.  On $E$, $\ad H$ is $-D$, so it has no
zero weight.  With $E_+$ denoting its positive-weight subspace, the
trace identity is
\begin{equation}
 \label{sym:algebra-trace}
 c\xi(p)=\tr(p|_{E_+}),\qquad
 p\in\mathfrak h_0,\qquad c=-h\ne0.
\end{equation}
Here $p|_{E_+}$ means the isotropy action by brackets.  In particular
$E_+$ corresponds to the geometrically stable plane.
\end{proposition}

\begin{proof}
Transitivity and the Hamiltonian bracket formula give
\[
 \xi(\{p,q\})=\omega_0(Y_p(0),Y_q(0)).
\]
Its radical is exactly the kernel of tangent evaluation.  This is
both the asserted symplectic identification and the assertion
$\mathfrak k=\operatorname{stab}_{\mathfrak h}\xi$.

Suppose an ideal $I$ lies in $\mathfrak k$ and $p\in I$.  Every
iterated bracket of its projected field with fields from
$\mathfrak h$ vanishes at zero.  If $Y_p$ had a first nonzero
homogeneous term of degree $j\ge1$, brackets with symmetries having
prescribed constant tangent values would differentiate that term.
After $j$ such brackets a nonzero constant value would result.
At each step the other bracket term has higher order, so it cannot
cancel this leading derivative.  This contradicts $I\subset
\mathfrak k$.  Hence $Y_p=0$ formally and $p$ is constant.  A constant
Hamiltonian has weight zero; Lemma~\ref{sym:trace} then gives $hp=0$.
Thus $I=0$.

For $p\in\mathfrak k$ and any $q\in\mathfrak h$, evaluation of the
projected bracket gives
$Y_{\{p,q\}}(0)=-dY_p(0)Y_q(0)$.  Apply this first to $p=H$ to obtain
the action $-D$, and then use Lemma~\ref{sym:trace} on the stable
plane to obtain Equation~\eqref{sym:algebra-trace}.
\end{proof}

\subsection{Homogeneous fields and their transport}

\begin{proposition}[Homogeneous formal fields]
\label{sym:homogeneous-fields}
Every $\mathfrak k$-invariant subspace $V\subset E$ has a unique
$\mathfrak h$-invariant formal distribution with value $V$ at zero.
The construction respects real structures, inclusions, and natural
tensor operations.  If $V=\mathfrak l/\mathfrak k$ for a subalgebra
$\mathfrak l\supset\mathfrak k$, this distribution is involutive.
\end{proposition}

\begin{proof}
Choose symmetries $p_1,\ldots,p_{2m}$ whose projected values form a
basis.  Write $G_t$ for the composition of their formal flows with
parameters $t_1,\ldots,t_{2m}$, and $F(t)=G_t(0)$.  Its linear term
is invertible, so the formal inverse function theorem gives $t=t(q)$.
Transport $V$ by $dG_t(0)$ and substitute $t(q)$.  All operations are
well-defined in formal displacement parameters of positive order.

We verify independence and invariance.  Two such parameterized
translations reaching the same point differ by a formal family
fixing zero.  Its right logarithmic derivative belongs
coefficientwise to the projected symmetry algebra: composition and
inversion give sums of conjugated generators, and their formal
adjoint series are iterated brackets.  Since the family fixes zero,
these derivatives belong to the projected isotropy.  Its derivative
at zero consequently preserves $V$, by the linear differential
equation starting at the identity.  This proves independence.
Including the formal flow of any further symmetry in the comparison
proves invariance.  Conversely any invariant field must be
transported from its value by these spanning flows, which proves
uniqueness.  The same argument applies to invariant tensors and to
the asserted compatibility operations.

For involutivity, the Levi bracket of an invariant distribution at
zero, evaluated on the values of $p,q\in\mathfrak l$, is the negative
of the projected bracket $\{p,q\}$ modulo $V$.  One can check this
without assuming that the fundamental fields themselves are tangent
everywhere: choose tangent sections with the same initial values
and use that their brackets with symmetry fields remain tangent.
Subtracting the two identities gives the claimed Levi bracket.
It vanishes since $\mathfrak l$ is a subalgebra.  Invariance and the
spanning formal translations make the obstruction vanish as a full
formal tensor.
\end{proof}

For the next proposition, an \emph{intrinsic} construction means one
that commutes with formal strict contact isomorphisms of the affine
data, including reciprocal constant rescalings of the two affine
fields.  Such rescalings do not change the symmetry equations.

\begin{proposition}[Intrinsic transport on an actual axis]
\label{sym:axis-transport}
On either entire typical strong axis there are local smooth formal
identifications of $(\mathfrak h,\mathfrak k,\xi)$, induced by strict
contact coordinate changes.  They may be chosen smoothly for a
smooth finite-dimensional family of paths on an axis patch.
Every intrinsic isotropy-invariant tensor or subspace construction,
extended homogeneously as in
Proposition~\ref{sym:homogeneous-fields}, has smooth full ambient
formal jets on that axis, and these jets are holonomic there.
The same conclusion holds for a continuous choice from a finite
intrinsic list of invariant subspaces after these identifications.
Natural choices commute with the dynamics.  Flow translation gives
the corresponding weak-axis statement, with almost-everywhere
agreement at regular centers.
\end{proposition}

\begin{proof}
Work first on a smooth axis patch in a fixed section, and let $a(t)$
be a path there.  By Proposition~\ref{sym:transitivity} and
Lemma~\ref{sym:full-projections}, the full finite symmetry projections
surject smoothly onto prescribed tangent values.  Choose a smooth
section with value $a'(t)$ at the first required order.  Lift it
successively through the surjective smooth bundle maps of
Lemma~\ref{sym:full-projections}, using smooth right inverses
chosen from bundle metrics on the axis.
This constructs a compatible family of formal symmetries $p_t$
whose every coefficient is smooth and for which
\[
 Y_{p_t}(a(t))=a'(t).
\]
The choice can be made smoothly in auxiliary path parameters by the
same finite-dimensional right inverses.  There is no assertion of
a bound uniform in the formal order.

Integrate these symmetries in coordinates relative to the moving
center.  In the section the generating field is
\[
 Y_{p_t}(a(t)+z)-a'(t),
\]
which vanishes at $z=0$.  Its ordinary differential equation therefore
closes at every finite Taylor order: the order-$N$ transported jet
depends only on the order-$N$ generator jet and lower orders.
Solving these compatible finite-dimensional equations produces a
smooth formal coordinate change.  For the strict contact lift,
choose the center height $\beta(t)$ by
\[
 \beta'(t)=p_t(a(t))-\theta_{a(t)}(a'(t)).
\]
The lifted generator relative to $(a(t),\beta(t))$ then also vanishes
at its formal origin.  The resulting formal changes are strict.

They intertwine the augmented affine fields.  Indeed differentiate
the pullback of their frame jets, modulo the corresponding bundle.
Moving-center differentiation is
the formal derivative in the center velocity, by whole-axis
holonomicity and the prescribed smooth height dependence.  The
relative generator subtracts that same velocity.  The remaining
term is its infinitesimal symmetry equation
\eqref{sym:linear-equations}, and is zero in the bundle quotient.
An internal change of frame is permitted.  This verification is made
at each finite order, where it is ordinary differentiation of a
finite system.  At the endpoint, the center height $\beta$ introduces
the reciprocal constants $e^{-h\beta}$ and $e^{h\beta}$ in the two
affine fields.  They do not change $\mathfrak h$.  Strictness and
transport of the center preserve both isotropy and the evaluation
functional $\xi$.

An intrinsic construction at the centers is thus constant under
these algebra identifications.  Proposition~\ref{sym:homogeneous-fields}
commutes with the identifications, so its entire formal field is
transported and has smooth coefficients.  To verify holonomicity,
let $T_t$ denote its moving-center formal tensor, projected to the
section.  For a distribution the following identities are interpreted
in a Grassmann chart.  The transport identity gives
\[
 \partial_t T_t+\mathcal L_{Y_{p_t}-a'(t)}T_t=0.
\]
Homogeneous invariance gives $\mathcal L_{Y_{p_t}}T_t=0$; hence
\[
 \partial_tT_t=\mathcal L_{a'(t)}T_t.
\]
For the constant displacement vector $a'(t)$ this is exactly the
coefficientwise Taylor coherence identity.  The argument for lifted
tensors includes the center height velocity in the same way.

For a finite intrinsic list, use the smooth family of paths to
trivialize the structures over a whole axis patch.  A continuous
selection into that finite list is locally constant in this
trivialization, and the preceding argument applies.  This patch
formulation also proves the assertion along paths that need not
themselves consist of regular points.  The whole-axis extensions of
Lemma~\ref{sym:full-projections} permit these patches everywhere on
the axis; uniqueness and regular density make their jets agree on
overlaps.  Naturality proves equivariance under the return map and
under regular flow shifts.  The prescribed height dependence,
strong-axis coherence, and flow Fubini then give the weak-axis
extension and its almost-everywhere agreement.
\end{proof}
